\lstset{
basicstyle=\fontsize{6.1}{8}\ttfamily,
    label={lst:train},
    numbers=left,
    caption={Prompting technique example for a loan application process (Bpi12), for the KPI of Total Time. Lines that have to be provided by the process analyst are marked in bold. In the example, only 2 training traces are provided due to space limitation.},
    captionpos=b,
    moredelim=**[is][\bfseries]{@}{@},
    float=t
} % Adjust scale as needed
\begin{lstlisting}
You are an expert in process mining and machine learning. Your task is to predict the 'total time' of 
process instances based on event logs, as each process instance is a sequence of activities.

A event log is a collection of traces, where each trace represents a process instance. 
Each trace is mapped as a sequence of activities and integers representing the minutes since the start
of the process.
The log is represented as a python list containing one dictionary for each trace. Included in it are:
@- the key "AMOUNT_REQ", representing the total amount of euros requested in the loan application.@
- the key "ActTimeSeq", which value is a list of [activity, cumulative elapsed minutes] 
- The key "total_time", which value is the total execution time in minutes from the start of the activity, 
that is the value to predict.

All interactions will be structured in the following way, with the appropriate values filled in.

[[ ## reasoning ## ]]
{your step-by-step reasoning}

[[ ## answer ## ]]
{your predicted total time as an integer}

[[ ## completed ## ]]

In adhering to this structure, your objective is to analyze the event log, and apply reasoning to predict 
the total time for a new case. This case belongs to a not-yet-completed process instance, represented by the 
label "Running" in "ActTimeSeq", indicating that more activities are expected before reaching the conclusion
of the process instance.

Ensure to articulate each step of your thought process in the reasoning field, detailing how you identify
relationships with past cases and leverage your intuition about the meaning of activities to arrive at the 
solution. The answer should be the final prediction of the total time for the given process instance.
Respond with the corresponding output fields, starting with the field [[ ## reasoning ## ]], 
then [[ ## answer ## ]], and then ending with the marker for [[ ## completed ## ]].

Your task is to learn from them and predict the 'total time' values for that traces.

@The process deals with a loan application process from a Dutch financial institution. It has been provided 
in the Business Process Intelligence (BPI) challenge in 2012. @
The following list shows some completed example cases with their total times:

     {"AMOUNT_REQ": 5000.0, "ActTimeSeq": [["W_Completeren aanvraag", 11],
     ["W_Nabellen offertes", 1464], ["W_Nabellen offertes", 7486]], "total_time": "7486"}
     {"AMOUNT_REQ": 15000.0, "ActTimeSeq": [["W_Completeren aanvraag", 13],
     ["W_Nabellen offertes", 14], ["W_Validate application", 4328], ["W_Validate application", 8792]], 
     "total_time": "8792"}
     ...
     ...
        
Now predict the total time for this new uncompleted case, considering that the case is still running:

    {"Application_1000386745": {"AMOUNT_REQ": 18000.0, 
    "ActTimeSeq": [["W_Completeren aanvraag", 2], ["W_Nabellen offertes", 8571], ["Running"]]}
\end{lstlisting}